\begin{slide}
\begin{center}\Title{Overview}\end{center}

\begin{itemize}

\IncPart{1}{\item[]
A function call in untyped $\lambda$-calculus corresponds to
a $\beta$-reduction step, whose most common form is: 
$\Comp{\Abst{x}T}\Appl{V} \BStep T\Subs{V}{x}$.
}

\IncPart{1}{\item[]
The substitution $T\Subs{V}{x}$
replaces with $V$ \stress{all} free instances of $x$ in $T$,
while avoiding captures through $\alpha$-conversion,
\ie, renaming.
}

\IncPart{2}{\item[]
From the computational standpoint,
avoiding captures is a time consuming operation
that reduction machines strive to avert.
}

\IncPart{3}{\item[]
Replacing named variable instances with de Bruijn's depth indexes
(\stress{nameless} binding)
speeds up this operation without nullifying it.
}

\IncPart{4}{\item[]
The depth index of a variable instance $x$ bound in $T$
is the number of $\lambda$'s of $T$ in whose scope $x$ occurs free.
Thus, it's a positive integer.
}

\IncPart{5}{\item[]
The nameless substitution $T\Subs{V}{1}$
may update the depth indexes of the free variables
in the copies of $V$ that it inserts in $T$.
}

\IncPart{6}{\item[]
This operation, known as lifting,
amounts to applying an \stress{update function} $\Uniform{h}$ (where $h \ge 0$)
to the depth indexes occurring in $V$.
}

\end{itemize}

\end{slide}
